% ================================================================
% PRIMARI
% ================================================================
\section{Processi Primari}

\subsection{Processo Di Fornitura}
\label{sec:fornitura}

\subsubsection{Descrizione}

Il processo di fornitura regola le attività con cui \textit{Synergo Unit} si pone
nella posizione di fornitore nei confronti del committente (Prof.\ Vardanega) e del
mentore/cliente (Zucchetti S.p.A.). Il Regolamento del Progetto Didattico stabilisce
le seguenti distinzioni di ruolo\textsubscript{G}:

\begin{itemize}
  \item il \textbf{proponente} è \textit{cliente} rispetto alle esigenze di prodotto
        e \textit{mentore} rispetto alle scelte di sviluppo;
  \item il \textbf{docente} è \textit{committente} di tutte le forniture;
  \item l'interazione fornitore-committente concerne vincoli contrattuali
        (scadenze, costi) e questioni regolamentari.
\end{itemize}

Attualmente il processo comprende la valutazione e la selezione del capitolato,
la comunicazione formale con i referenti esterni e la produzione della documentazione
richiesta per la candidatura\textsubscript{G}.

\subsubsection{Attività}

\paragraph{Valutazione Dei Capitolati\textsubscript{G}}
Il gruppo analizza i capitolati disponibili producendo il documento
\href{https://synergo-unit-unipd.github.io/Documentary-Repo/docs/CC/doc_gestionali/doc_esterni/valutazione_dei_capitolati/valutazione_dei_capitolati.pdf}
     {\underline{Valutazione Dei Capitolati}},
che illustra i criteri di confronto, i vantaggi e gli svantaggi di ciascun
capitolato — includendo il resoconto degli incontri esplorativi con i proponenti —
e motiva la scelta finale. Il documento è soggetto alla revisione\textsubscript{G}
interna prima della pubblicazione.

\paragraph{Dichiarazione Degli Impegni\textsubscript{G}}
A seguito della selezione del capitolato, il gruppo redige la
\href{https://synergo-unit-unipd.github.io/Documentary-Repo/docs/CC/doc_gestionali/doc_esterni/dichiarazione_degli_impegni/dichiarazione_degli_impegni.pdf}
     {\underline{Dichiarazione Degli Impegni}},
che specifica:

\begin{itemize}
  \item scadenza ultima di consegna prevista;
  \item preventivo dei costi\textsubscript{G} finali calcolato secondo il
        regolamento (costo minimo ammesso: 12.000\,€ per gruppi di 7,
        proporzionalmente ridotto per gruppi di 6);
  \item totale delle ore produttive assegnate a persona, strettamente compreso
        nell'intervallo 80\,(min)~--~95\,(max) ore per membro;
  \item distribuzione delle ore per ruolo (non per persona) secondo i costi orari
        ufficiali riportati nella Sezione~\ref{sec:ruoli};
  \item rischi attesi e strategie di mitigazione;
  \item impegni formali assunti verso il committente.
\end{itemize}

Le ore produttive sono contabilizzate solo dopo l'aggiudicazione del capitolato.
Il costo accettato all'aggiudicazione diventa limite superiore invalicabile durante
lo svolgimento del progetto.

\paragraph{Comunicazioni Esterne}
Le comunicazioni ufficiali con aziende e docenti avvengono tramite l'indirizzo email
istituzionale del gruppo. Tale indirizzo è configurato con inoltro automatico verso
gli indirizzi personali dei membri delegati; le risposte vengono inviate mantenendo
come mittente l'indirizzo ufficiale, garantendo continuità e tracciabilità delle
comunicazioni formali. Le candidature alle revisioni di avanzamento avvengono per
posta elettronica, senza invio di allegati, esponendo il puntatore alla sezione del
repository\textsubscript{G} dedicata alla revisione.

\paragraph{Incontri Con Il Proponente}
Ogni incontro con il proponente o con i docenti viene pianificato con congruo
anticipo, prevede la predisposizione di un ordine del giorno e si conclude con la
redazione e approvazione di un verbale esterno secondo le norme descritte nella
Sezione~\ref{sec:documentazione}. La frequenza degli incontri col proponente/mentore
deve essere la più alta possibile.

\paragraph{Diario Di Bordo}
Il gruppo è tenuto a relazionare al committente negli eventi diario di bordo
secondo il calendario stabilito, esponendo il rapporto tra costi sostenuti e
avanzamento conseguito.